\chapter{Digital-Twin Architecture}
\label{chap:architecture}

This chapter describes the design and implementation of the discrete-event
simulation framework that forms the computational core of the DDP.  Each
section maps to one or more modules in the \texttt{engine/} package and
explains the algorithmic choices made at that layer.

%% -----------------------------------------------------------------------
\section{Network Topology}
\label{sec:arch:topology}

The framework represents any supply-chain network as a \textbf{directed
acyclic graph (DAG)}, implemented in \texttt{engine/network.py}.  Each
\emph{node} in the graph corresponds to a stocking location (supplier,
warehouse, or retailer) and each directed \emph{edge} to a transport lane
between two locations.  Goods flow strictly downstream; returns and lateral
transshipment are not modelled.

The topology is entirely configurable: the number of echelons (chain depth)
and the number of nodes at each echelon (chain breadth) are controlled by the
JSON configuration file without any change to simulator code.  A single
supplier feeding one warehouse and three retailers, a two-DC hub-and-spoke
network, or a flat one-level distribution system are all expressible as DAG
configurations.  Throughout this work, several topologies of varying size and
depth were used; the specific configuration for each experiment is described
in Chapter~\ref{chap:methodology}.

Two further design choices are worth noting.  First, the supplier node carries
an optional \texttt{infinite\_supply} flag: when set, the supplier can always
fulfil any order placed by downstream nodes, representing an ideally reliable
upstream source.  This assumption is convenient for isolating the behaviour
of downstream inventory policies but is entirely optional — finite-capacity
suppliers and supply disruptions are also supported
(Section~\ref{sec:arch:disruptions}).  Second, the DAG structure enforces a
strict material-flow direction, ensuring that cycle detection is trivial and
the topological processing order (Section~\ref{sec:arch:des}) is unique.

\begin{figure}[htbp]
  \centering
  \fbox{\parbox{0.65\textwidth}{\centering
    \textit{[Figure: Example multi-echelon supply network — Supplier
    $\to$ Warehouses W1, W2 $\to$ Retailers R1, R2, R3.
    Arrows show downstream material flow; lead times labelled on each lane.]}}}
  \caption{Example DAG topology: one supplier, two warehouses, and three
           retailers.  The framework supports arbitrary depth and breadth;
           this is one of several topologies used during development.}
  \label{fig:topology}
\end{figure}

%% -----------------------------------------------------------------------
\section{Discrete-Event Simulation Loop}
\label{sec:arch:des}

The simulator advances time in unit steps of one day.  Each day $t$ executes
the following sequence of operations:

\begin{enumerate}
  \item \textbf{Receive shipments.}  For every node, any shipment whose
        scheduled arrival time equals $t$ is removed from the pipeline and
        added to on-hand inventory.
  \item \textbf{Realise demand.}  Each retailer node draws its external demand
        for day $t$ from the configured demand generator (deterministic,
        Poisson, or CSV-driven).
  \item \textbf{Clear external backlog.}  Any accumulated external demand that
        could not be served on previous days is fulfilled from current on-hand
        stock, up to availability.
  \item \textbf{Process child orders.}  Internal orders from downstream nodes
        are served from the parent's on-hand inventory; unmet orders enter the
        parent's \emph{children backlog}.
  \item \textbf{Place replenishment orders.}  Each non-supplier node queries its
        inventory policy to determine an order quantity and registers an
        \texttt{IncomingOrder} at its parent.
  \item \textbf{Dispatch planning.}  For each edge, the transport module
        inspects accumulated orders, builds vehicle loads respecting capacity
        constraints, and schedules shipments with the configured lead time
        (Section~\ref{sec:arch:transport}).
  \item \textbf{Cost accumulation.}  Holding, shortage, ordering, and
        transport costs for day $t$ are recorded.
\end{enumerate}

Nodes are processed in \emph{topological order} — retailers before warehouses,
warehouse before the supplier — so that each node's child-order step is
complete before its parent's dispatch step.  The topological sort is computed
once at simulation startup (\texttt{Simulator.\_topological\_order}).

The entry point is \texttt{Simulator.run(mode)}, implemented in
\texttt{engine/simulator.py}.  When \texttt{mode="logs"} the simulator records
full daily time-series; when \texttt{mode="summary"} only the aggregate KPI
row is returned, which is sufficient for the optimisation inner loop.

\subsection{Pseudocode}

Algorithm~\ref{alg:simloop} summarises one daily tick.

\begin{algorithm}
\caption{Daily simulation tick at time $t$}
\label{alg:simloop}
\begin{algorithmic}[1]
\For{each node $n$ in topological order}
  \State $n.\text{receive\_shipments}(t)$
  \If{$n$ is a retailer}
    \State demand $\leftarrow$ demandGenerator$[n].\text{get}(t)$
    \State $n.\text{process\_external\_demand}(\text{demand}, t)$
  \EndIf
  \State $n.\text{clear\_external\_backlog}()$
  \State $n.\text{process\_child\_orders}(t)$
\EndFor
\For{each non-supplier node $n$ in topological order}
  \State $q \leftarrow n.\text{policy}.\text{order\_qty}(
         \text{on\_hand}, \text{backlog}, \text{pipeline\_in}, t)$
         \Comment{$q$: order quantity in units}
  \State $n.\text{parent}.\text{add\_inbound\_order}(n, q, t)$
\EndFor
\For{each edge $e$ in the network}
  \State loads $\leftarrow$ GreedyLoadPlanner$.\text{plan}(e.\text{pending\_orders})$
  \For{each load $\ell$ in loads}
    \State schedule shipment arrival at $t + e.\text{lead\_time}$
    \State accumulateCost$(e, \ell)$
  \EndFor
\EndFor
\State accumulateCost$(t)$ \Comment{holding and shortage}
\end{algorithmic}
\end{algorithm}

%% -----------------------------------------------------------------------
\section{Node State Model}
\label{sec:arch:node}

Each node in the network maintains the following per-SKU state variables,
implemented as dictionaries keyed by SKU identifier in \texttt{engine/node.py}:

\begin{description}
  \item[\texttt{on\_hand}] Current physical stock available for dispatch or
        external fulfilment.
  \item[\texttt{pipeline\_in}] A list of \texttt{Shipment} objects, each
        carrying a quantity and a scheduled arrival day.  The total pipeline
        quantity is the sum of all in-transit units.
  \item[\texttt{backlog\_external}] Unfulfilled external demand accumulated
        over previous days; cleared before new demand is processed on each
        tick.
  \item[\texttt{backlog\_children}] A dictionary keyed by \texttt{(child\_id,
        sku)} recording how many units the parent owes to each downstream
        child.
  \item[\texttt{inbound\_orders\_queue}] Orders submitted by children that
        have not yet been dispatched.
  \item[\texttt{placed\_orders}] Historical record of orders placed by this
        node; used to compute the bullwhip ratio.
\end{description}

The two backlog dictionaries encode the dual role of intermediate nodes:
they both \emph{sell} to their children (backlog\_children) and
\emph{buy} from their parent (backlog\_external records their own unmet need).

\subsection{Helper methods}

\texttt{Node.total\_pipeline\_in(sku)} returns the scalar sum over all
in-transit shipments for a given SKU, which is the quantity passed to the
inventory policy as \texttt{pipeline\_in}.

\texttt{Node.total\_backlog\_children(sku)} aggregates children-owed
quantities across all child nodes, used by the echelon-stock calculation
(Section~\ref{sec:policies:echelon}).

%% -----------------------------------------------------------------------
\section{Transport Layer}
\label{sec:arch:transport}

Transport between nodes is modelled at the \emph{vehicle load} level rather
than at the individual unit level.  The transport layer is implemented in
\texttt{engine/transport.py} and consists of three collaborating types.

A key design feature is that a single lane (edge) between two nodes can carry
\textbf{multiple transport options} simultaneously — one per available mode.
For example, a warehouse-to-retailer lane might offer both a truck option
(high capacity, 1-day lead time) and a rail option (higher capacity, 3-day
lead time, lower cost per vehicle).  This directly mirrors real-world
logistics, where two locations are often served by more than one carrier type,
each with its own cost structure and transit time.  Multiple options are
declared by listing several edge entries with the same \texttt{from}/\texttt{to}
pair in the config's \texttt{edges} array, each with a distinct \texttt{mode}
integer.

\subsection{TransportOption}

A \texttt{TransportOption} specifies the physical and economic properties of
one transport mode on a lane:

\begin{itemize}
  \item \textbf{mode} — a string label (e.g.\ \texttt{"truck"}, \texttt{"rail"}).
  \item \textbf{capacity} — maximum volume (in cubic metres) per vehicle.
  \item \textbf{weight\_capacity} — maximum gross weight (kg) per vehicle;
        \texttt{None} means no weight constraint.
  \item \textbf{cost\_full / cost\_half / cost\_quarter} — tiered cost
        schedule: a full load costs \texttt{cost\_full}, a load between 50\%
        and 100\% utilisation costs \texttt{cost\_half}, and a load below 50\%
        costs \texttt{cost\_quarter}.
  \item \textbf{lead\_time} — deterministic transit days (stochastic lead times
        are a planned extension).
  \item \textbf{min\_dispatch\_utilization} — minimum fraction of capacity
        that must be filled before a vehicle is released.  A value of 0 means
        dispatch on any positive load; a value of 0.8 means the planner
        waits until 80\% of capacity is committed.
\end{itemize}

When multiple \texttt{TransportOption}s exist on a lane, the framework's
default \textbf{mode selector} sorts them by their \texttt{mode} integer
(1~=~primary, 2~=~secondary, \ldots) and dispatches all cargo via the
lowest-numbered option, with higher-numbered modes serving as fallbacks.
In addition, the optional \texttt{share} field on an edge enables
\emph{weighted random sampling} of lead times across options — for example,
assigning \texttt{share: 0.8} to a truck route and \texttt{share: 0.2} to
a rail route means 80\% of dispatch events use the truck lead time and 20\%
use the rail lead time, directly modelling realistic carrier variability for
multi-route logistics corridors without splitting individual vehicle loads.

\subsection{GreedyLoadPlanner}

The \texttt{GreedyLoadPlanner} is the load-packing algorithm used in all
experiments in this work.  It is the default implementation of the
\texttt{LoadPlanner} protocol.  Given a list of pending order lines for a
lane, it greedily packs items into vehicles one at a time, subject to dual
constraints:

\begin{enumerate}
  \item \textbf{Volume constraint:} total volume of packed items
        $\leq$ \texttt{capacity}.
  \item \textbf{Weight constraint:} total weight of packed items
        $\leq$ \texttt{weight\_capacity} (when set).
\end{enumerate}

If weight capacity is tighter than volume capacity for the current mix of SKUs,
the effective capacity is scaled down accordingly before packing begins.  After
packing, if the resulting utilisation is below \texttt{min\_dispatch\_utilization},
the load is \emph{held} for up to \texttt{max\_dispatch\_wait} days; once the
wait limit is reached the load is dispatched regardless of fill level to prevent
indefinite stockouts at downstream nodes.

The tiered cost is evaluated after packing: loads $\geq 50\%$ full pay the
half-load rate, loads $\geq 100\%$ (i.e.\ exactly full) pay the full-load rate.
This incentivises consolidation but never blocks dispatch beyond the wait limit.

\subsection{Pluggable protocols}

The transport layer is designed around three narrow protocols, making all
three components fully \textbf{plug-and-play}:

\begin{description}
  \item[\texttt{LoadPlanner}] Decides how pending orders are packed into
        vehicle loads.  \texttt{GreedyLoadPlanner} is the default used in
        this work; alternatives (bin-packing, probabilistic packing) can be
        substituted by implementing the same two-method interface.
  \item[\texttt{ModeSelector}] Chooses which transport option(s) to activate
        when multiple options exist on a lane.  The default
        \texttt{CheapestModeSelector} always picks the primary mode
        (lowest \texttt{mode} integer); a custom selector could split loads
        across modes or apply routing rules.
  \item[\texttt{CostAllocator}] Distributes a vehicle's total cost across the
        SKUs it carries.  The default allocates cost proportionally by volume.
\end{description}

No changes to the simulator core (\texttt{engine/simulator.py}) are required
to swap any of these components — they are injected at construction time,
making the transport layer independently extensible.

%% -----------------------------------------------------------------------
\section{Cost Model}
\label{sec:arch:cost}

The total cost accumulated over the simulation horizon $T$ is:

\begin{equation}
  C_{\text{total}} = C_{\text{hold}} + C_{\text{transport}}
                   + C_{\text{order}} + C_{\text{shortage}}
  \label{eq:total_cost}
\end{equation}

Each component is defined as follows.

\paragraph{Holding cost.}
$C_{\text{hold}} = \sum_{t=w}^{T} \sum_{n} \sum_{s}
  h_{n,s} \cdot \text{on\_hand}_{n,s,t}$,
where $h_{n,s}$ is the per-unit-per-day holding cost for SKU $s$ at node $n$
and $w$ is the warm-up cutoff.

\paragraph{Shortage (backlog) cost.}
$C_{\text{shortage}} = \sum_{t=w}^{T} \sum_{n} \sum_{s}
  p_{n,s} \cdot \text{backlog\_external}_{n,s,t}$,
where $p_{n,s}$ is the per-unit-per-day penalty cost.

\paragraph{Ordering cost.}
Each order placed by a node incurs a fixed setup cost $K_n$ plus a
per-unit variable cost $c_{n,s}$:
$C_{\text{order}} = \sum_{\text{orders}} (K_n + c_{n,s} \cdot q)$,
where $q$ is the order quantity in units.

\paragraph{Transport cost.}
Each dispatched vehicle load on edge $e$ with utilisation $u$ incurs
a tiered cost: $\text{cost\_full}$ if $u = 1$, $\text{cost\_half}$ if
$0.5 \leq u < 1$, and $\text{cost\_quarter}$ if $u < 0.5$.

All four components are accumulated on every tick inside
\texttt{Simulator.run} and the first $w$ days are excluded from reported
KPIs to avoid contamination by the initial-inventory transient
(Section~\ref{sec:meth:warmup}).

%% -----------------------------------------------------------------------
\section{Configuration Schema}
\label{sec:arch:config}

Scenarios are specified as JSON files.  The top-level keys are:

\begin{description}
  \item[\texttt{seed}] Integer random seed for reproducibility.
  \item[\texttt{time\_horizon}] Number of simulation days.
  \item[\texttt{skus}] List of SKU identifiers.
  \item[\texttt{sku\_properties}] Per-SKU \texttt{volume\_per\_unit} and
        \texttt{weight\_per\_unit}.
  \item[\texttt{nodes}] List of node objects specifying type
        (\texttt{supplier}/\texttt{warehouse}/\texttt{retailer}),
        initial inventory, cost parameters, and per-SKU policy
        configuration.
  \item[\texttt{edges}] List of edge objects specifying from/to nodes and
        a list of \texttt{TransportOption} specifications.
  \item[\texttt{disruptions}] (Optional) List of supply disruption events;
        each event specifies a node, a start day, and a duration.
\end{description}

The \texttt{engine/config\_validator.py} module validates all required fields
and raises descriptive errors before the simulation starts, preventing
silent misconfiguration during long experiment runs.

Randomness in the framework arises from two sources: (i)~stochastic demand
generators — for example, Poisson-distributed daily demand draws a fresh
random variate each day — and (ii)~the Optuna TPE optimiser, which uses
probabilistic surrogate models and random sampling to explore the parameter
space.  The \texttt{seed} value initialises Python's \texttt{random} module
and NumPy's RNG at simulation startup, ensuring that both demand sequences
and optimiser trial sequences are exactly reproducible across runs.  CSV-driven
demand (used in all main experiments of this work) is fully deterministic and
unaffected by the seed; the seed matters primarily when using Poisson or
other synthetic demand modes and when running the optimiser.

%% -----------------------------------------------------------------------
\section{CSV Input Pipeline}
\label{sec:arch:csv}

To lower the barrier for practitioners who are not comfortable editing JSON
directly, a secondary input method was developed using six CSV files.  The
builder script \texttt{inputs/scripts/build\_config\_from\_csv.py} reads and
validates all six files and emits a complete JSON configuration that can then
be passed to any simulation script.

Table~\ref{tab:csv-files} lists the six files and their purposes.

\begin{table}[htbp]
  \centering
  \caption{CSV input files and their roles in the configuration pipeline.}
  \label{tab:csv-files}
  \begin{tabular}{@{}ll@{}}
    \toprule
    File & Contents \\
    \midrule
    \texttt{sim\_config.csv}    & Global parameters: seed, horizon, warm-up, strategy \\
    \texttt{products.csv}       & SKU physical properties: volume and unit cost \\
    \texttt{nodes.csv}          & One row per node — type, inventory, cost parameters \\
    \texttt{policies.csv}       & Long-format policy parameters (one row per parameter) \\
    \texttt{routes.csv}         & Transport lanes: capacity, tiered costs, lead time \\
    \texttt{demand\_config.csv} & Maps each retailer to its demand CSV \\
    \bottomrule
  \end{tabular}
\end{table}

Per-retailer daily demand is read from CSV files in \texttt{inputs/demand\_data/},
derived from the M5 Walmart Forecasting dataset.

%% -----------------------------------------------------------------------
\section{Supply Disruption Modelling}
\label{sec:arch:disruptions}

Supply disruptions are modelled by temporarily blocking a node from
receiving or dispatching stock for a specified interval $[t_s, t_s+d)$.
During a disruption, orders placed by downstream nodes accumulate in the
\texttt{backlog\_children} dictionary of the affected node.  When the node
recovers at day $t_s + d$, it resumes normal order fulfilment and works
through the accumulated backlog.

Disruptions are declared in the \texttt{disruptions} list of the config JSON.
Multiple, non-overlapping disruption events can be specified per node.  The
disruption model is deliberately simple — it does not include rerouting or
mode switching — making it straightforward to extend.

%% -----------------------------------------------------------------------
\section{Iterative Framework Evolution}
\label{sec:arch:evolution}

The simulator was built and validated incrementally across three phases.

\textbf{Phase 1 — Prototype and validation.}
The foundation of the framework was built as a single-SKU discrete-event
engine and tested across several small topologies: 1-1-1, 1-1-2, 1-2-3,
and other configurations.  Both synthetic demand profiles (constant demand
and Poisson-distributed daily draws) and real M5 Walmart CSV demand were
used as inputs.  The goal was to ensure that all core modules — the
daily-tick loop, shipment pipeline, backlog propagation, and policy
interface — behave correctly and agree with theoretical expectations.
This included verifying lead-time mechanics (shipments placed on day $t$
arrive on day $t+L$ with the correct quantity) alongside all other
invariants.  Four controlled validation experiments (EV1--EV4,
Chapter~\ref{chap:methodology}) formalised this verification before any
additional complexity was introduced.

\textbf{Phase 2 — Transport layer, CSV pipeline, and scale-up.}
The most significant architectural addition in this phase was the
\textbf{transport layer}: vehicle-capacity constraints (simultaneous volume
and weight limits), tiered cost schedules (full-, half-, and quarter-load
rates), configurable minimum-dispatch-utilisation thresholds, and a
maximum-dispatch-wait timeout to prevent indefinite consolidation delays.
This transformed the simulator from a pure inventory tool into a joint
inventory-transport digital twin.  A \textbf{CSV input pipeline} was also
introduced, providing a spreadsheet-friendly alternative to hand-editing
JSON configurations.  Concurrently, the network was extended to multi-SKU
scenarios, the policy library grew from two implementations (base-stock,
$(s,S)$) to seven, and multi-seed validation and bullwhip metrics were
instrumented.

\textbf{Phase 3 — Optimisation integration.}
The Optuna-TPE optimiser was integrated and the main experiments (E0--E4)
were designed and executed.  The joint optimiser searches inventory
base-stock levels and transport dispatch thresholds simultaneously — a
search structure only made meaningful by the transport layer built in
Phase~2.  Full details of the optimisation framework are given in
Chapter~\ref{chap:optimization}.

This incremental approach ensured each layer was tested in isolation before
the next layer of complexity was added, keeping the test suite tractable at
every stage.

\paragraph{Test suite.}
Throughout all three phases, a growing \texttt{pytest} suite was maintained
to guard the correctness of core mechanics.  Tests cover simulator invariants
(pipeline quantity conservation, backlog non-negativity, cost accumulation
bounds), policy behaviour under edge cases (zero demand, large backlogs,
review-day alignment), and transport layer mechanics (dual-constraint
packing, cost-tier selection, dispatch-wait timeout logic).  Running
\texttt{python -m pytest tests/ -q} before each new phase confirmed that
foundational logic remained intact as the framework grew in complexity,
and the suite continues to serve as a regression guard for future
extensions.
